% Generated draft: requires mathematical and source review.
\documentclass[11pt]{article}
\usepackage[margin=25mm]{geometry}
\usepackage{fontspec}
\setmainfont{Latin Modern Roman}
\usepackage{amsmath,amssymb,mathtools,unicode-math}
\setmathfont{Latin Modern Math}
\usepackage{enumitem,needspace,xurl,hyperref}
\setlength{\parindent}{0pt}
\setlength{\parskip}{5pt}
\setlength{\emergencystretch}{3em}
\newcommand{\N}{\mathbb{N}}
\newcommand{\Z}{\mathbb{Z}}
\newcommand{\Q}{\mathbb{Q}}
\newcommand{\R}{\mathbb{R}}
\newcommand{\C}{\mathbb{C}}
\newcommand{\F}{\mathbb{F}}
\newcommand{\A}{\mathbb{A}}
\newcommand{\PP}{\mathbb{P}}
\newcommand{\E}{\mathbb{E}}
\newcommand{\eps}{\varepsilon}
\newcommand{\dd}{\,\mathrm{d}}
\begin{document}
\section*{UnsolvedMath ID: 800360}
\textbf{Source problem code:} catalogue entry
\subsection{Definitions and mathematical statement}
Let \(I\) be a finite set, and let \((m_{ij})_{i,j\in I}\) be a symmetric matrix with \(m_{ii}=1\) and \(m_{ij}\in\{2,3,\ldots\}\cup\{\infty\}\) for \(i\ne j\). Define the Coxeter group
\[
W=\left\langle s_i\ (i\in I)\ \middle|\ (s_is_j)^{m_{ij}}=1\text{ whenever }m_{ij}<\infty\right\rangle.
\]
For \(w\in W\), let \(\ell(w)\) be the minimum number of generators needed to express \(w\). A positive weight function is a map \(L:W\to\mathbb Z_{\ge0}\) such that \(L(w)>0\) for \(w\ne1\) and
\[
\ell(ww')=\ell(w)+\ell(w')\quad\Longrightarrow\quad L(ww')=L(w)+L(w')
\qquad(w,w'\in W).
\]
These conditions imply \(L(1)=0\); thus the source's notation \(\mathbb N\) must here allow zero. The length function itself is an example.

Let \(v\) be an indeterminate and set \(A=\mathbb Z[v,v^{-1}]\). The Iwahori--Hecke algebra \(H=H(W,L)\) is the free \(A\)-module with basis \(\{T_w:w\in W\}\), equipped with the associative \(A\)-algebra structure determined by
\[
(T_{s_i}+v^{-L(s_i)})(T_{s_i}-v^{L(s_i)})=0\quad(i\in I),
\qquad
T_wT_{w'}=T_{ww'}\quad\text{if }\ell(ww')=\ell(w)+\ell(w').
\]
For \(c\in\mathbb C\setminus\{0\}\), define
\[
H_c=\mathbb C\otimes_{A,\,v\mapsto c}H.
\]
For an affine Weyl group \(W\), let \(\Sigma_W\) denote the collection of weight functions referred to in the source as being described in [L91, section 17], [L95], and [L02]. This is a reference-dependent designation: the supplied record does not give a general definition of \(\Sigma_W\), and the section reference to [L91] has damaged formatting.

\textbf{Requested assertion, with unresolved source qualifications.} For every algebra associated in [S20] to a supercuspidal representation of a parabolic subgroup of a \(p\)-adic reductive group, show that, up to extension by a group algebra of a small finite group, it is of the form \(H_q\), where \(q\) is a power of \(p\), \(W\) is an affine Weyl group, and \(L\in\Sigma_W\).

The preceding wording retains the source's phrase ``supercuspidal representation of a parabolic subgroup.'' The supplied record does not identify the precise algebras constructed in [S20], specify whether a representation of a Levi factor is intended, define the extension operation, or make ``small finite group'' quantitative. Consequently it does not determine a fully explicit algebra-isomorphism assertion without further source clarification.

The source gives the following examples. In affine type \(F_4\), the allowed tuples of values of \(L\) on the simple generators are exactly
\[
(1,1,1,1,1),\quad(1,1,1,2,2),\quad(2,2,2,1,1),\quad(1,1,1,4,4).
\]
In affine type \(G_2\), they are exactly
\[
(1,1,1),\quad(1,1,3),\quad(3,3,1),\quad(1,1,9).
\]
These lists use the source's implicit ordering of the simple generators; that ordering is not supplied in this record.

As background, the source states that the analogous assertion for groups with connected centre over a finite field \(\mathbb F_q\) is known, with the parenthetical allowance for a finite-group-algebra extension omitted. In that assertion, \(W\) is a Weyl group and \(\Sigma_W\) is the collection described in [L78, p.\ 35]. This is a supplied contextual claim, not an independently verified theorem or a replacement for the \(p\)-adic assertion.
\subsection{Short English statement}
Lusztig asks whether the algebras referred to in [S20], arising from supercuspidal representation data for \(p\)-adic reductive groups, have affine Iwahori--Hecke form \(H_q\) with a prescribed weight function, allowing a finite group algebra extension. The construction, the extension operation, and the general prescribed collection require clarification from the cited sources. The registry labels this problem open; this draft does not independently verify that status.
\subsection{Sources}
\begin{itemize}
\item Local UnsolvedMath registry: permanent ID \(800360\), exact source code \texttt{catalogue entry}, title ``Iwahori--Hecke Algebras --- Problem 1.'' Attribution: Lusztig; proposed year: 2020. The supplied registry status is \texttt{open}.
\item Lusztig, \emph{Open Problems on Iwahori-Hecke Algebras} (2020), Problem 1, \url{https://arxiv.org/abs/2006.08535}. The supplied provenance records access on 2026-10-16 and extraction by \texttt{arxiv\_latex}; this access information is supplied provenance and was not independently verified for this draft.
\item [S20], cited in the supplied problem as the source of the representation-theoretic algebras. No author, title, bibliographic details, or separate link for this reference is supplied.
\item [L91, section 17], [L95], and [L02], cited in the supplied problem for the affine collections \(\Sigma_W\). The section designation in [L91] has damaged source formatting. Full bibliographic entries and separate links are not supplied.
\item [L78, p.\ 35], cited in the supplied problem for the finite-field collections of weight functions. A full bibliographic entry and separate link are not supplied.
\end{itemize}
\end{document}
